% This is made in A4, so the text is really big when scaled up to A3.
\documentclass[20pt]{beamer}
\usetheme{default}
\usefonttheme{serif}
\setbeamertemplate{navigation symbols}{}
\geometry{a4paper,landscape}

\usepackage{minted}

\author{Michael Walker}
\title{A Sane Haskell Development Environment}
\institute{York Haskell\\
  \texttt{mike@barrucadu.co.uk}
}
\date{}

\begin{document}

\begin{frame}[plain]
  \titlepage
\end{frame}

\begin{frame}[fragile]{Hackage}
  \begin{verbatim}
msw504 on pc116 in /tmp/demo 
 >>>  cabal update
Downloading the latest package list from hackage.haskell.org

msw504 on pc116 in /tmp/demo 
 >>>  cabal install hello
Resolving dependencies...
Configuring hello-1.0.0.2...
Building hello-1.0.0.2...
Preprocessing executable 'hello' for hello-1.0.0.2...
[1 of 1] Compiling Main             ( src/hello.hs,
dist/dist-sandbox-d9586924/build/hello/hello-tmp/Main.o )
Linking dist/dist-sandbox-d9586924/build/hello/hello ...
Installing executable(s) in /tmp/demo/.cabal-sandbox/bin
Installed hello-1.0.0.2

msw504 on pc116 in /tmp/demo 
 >>>  ./.cabal-sandbox/bin/hello 
Hello, World!
  \end{verbatim}
\end{frame}

\begin{frame}[fragile]{Building with Cabal}
  \begin{columns}

  \column{0.34\textwidth}
  \texttt{ghc Main.hs}\\\texttt{ghc --make Main.hs}\\\texttt{runhaskell Main.hs}

  \vspace{3cm}

  \begin{enumerate}
  \item Make a .cabal file
  \item Fill in your dependencies
  \item \texttt{cabal build}
  \item PROFIT!
  \end{enumerate}

  \vspace{3.5cm}

  \column{0.66\textwidth}
  \begin{verbatim}
name:               demo
version:            0.1.0.0
synopsis:           This is a simple .cabal file
homepage:           http://www.example.com
license:            PublicDomain
license-file:       LICENSE
author:             Michael Walker
maintainer:         mike@barrucadu.co.uk
category:           Development
build-type:         Simple
cabal-version:      >=1.10

library
  exposed-modules:  Foo
  build-depends:    base >=4.6 && <4.7
  default-language: Haskell2010
  \end{verbatim}
  \end{columns}
\end{frame}

\begin{frame}[fragile]{Cabal Hell}
  \begin{columns}
  \column{0.5\textwidth}
  \begin{itemize}
  \item Hackage is not a curated collection of packages
  \item Package dependency constraints may be too strict
  \item Packages may have incompatible dependencies
  \item There is no guarantee that any two packages will work together
  \end{itemize}

  \column{0.5\textwidth}
  \begin{itemize}
  \item Stackage is a curated collection of packages
  \item Any two packages should work together
  \item Has fewer packages than Hackage
  \end{itemize}
  \vspace{5cm}
  \end{columns}

  \hrule

  \begin{verbatim}
msw504 on pc116 in /tmp/demo 
 >>>  cabal sandbox init
Writing a default package environment file to
/tmp/demo/cabal.sandbox.config
Creating a new sandbox at /tmp/demo/.cabal-sandbox

msw504 on pc116 in /tmp/demo 
 >>>  cabal install uniplate
Resolving dependencies...
 -- output omitted --
Installed uniplate-1.6.12
  \end{verbatim}
\end{frame}

\begin{frame}[fragile]{Documenting with Haddock}
  \begin{minted}{haskell}
-- | A value of type @Find m a@ represents a concurrent search
-- computation (happening in the 'MonadConc' monad @m@) which may
-- produce a value of type @a@, or fail. If a value can be returned,
-- one will be (although it's nondeterministic which one will actually
-- be returned). Usually you will be working with values of type @Find
-- IO a@, but the generality allows for testing.
--
-- You should prefer using the 'Applicative' instance over the 'Monad'
-- instance if you can, as the 'Applicative' preserves parallelism.
newtype Find m a = Find { unFind :: m (WorkItem m a) }
  \end{minted}

  \vspace{2cm}

  \begin{center}
    \url{http://barrucadu.github.io/search-party/}
  \end{center}
\end{frame}

\begin{frame}{Version Control}

  \begin{center}
    \includegraphics[width=4cm]{git.eps}
    \hspace{3cm}
    \includegraphics[width=4cm]{darcs.eps}
  \end{center}

  \vspace{2cm}

  A \textit{lot} of Haskell packages are on GitHub. This means that
  there's a lot to learn from!

  \vspace{5cm}
\end{frame}

\begin{frame}[fragile]{Continuous Integration}
\begin{verbatim}
language: haskell

env:
 - GHCVER=7.6.3
 - GHCVER=7.8.4

matrix:
  allow_failures:
   - env: GHCVER=head

before_install:
  ...

install:
  - $CABAL install --dependencies-only --enable-tests
  - $CABAL configure -flib-Werror --enable-tests $MODE

script:
 - $CABAL configure --enable-tests
 - $CABAL build
 - $CABAL test
\end{verbatim}
\end{frame}

\begin{frame}{My Workflow}
  \begin{enumerate}
  \item Agonise over name
  \item \texttt{cabal init}
  \item Ignore licensing as usual, and default to the WTFPL
  \item Hack away until something substantial exists
  \item Panic, and realise it's all untested
  \item Add a \texttt{.travis.yml} to build with different GHC
    versions
  \item Fix bugs
  \item Write tests
  \item Run tests in \texttt{.travis.yml}
  \item Fix more bugs
  \item Feel simultaneously guilty over taking so long to ``finish'',
    and inadequate to push to Hackage
  \item Continue!
  \end{enumerate}
\end{frame}

\begin{frame}{Resources}
  \begin{center}
  \url{https://www.haskell.org/cabal/users-guide/}

  \vspace{2cm}

  \url{https://www.haskell.org/haddock/doc/html/}

  \vspace{2cm}

  \url{https://downloads.haskell.org/~ghc/latest/docs/html/users\_guide/}

  \vspace{2cm}

  \url{http://eagain.net/articles/git-for-computer-scientists/}

  \vspace{2cm}

  \url{https://travis-ci.org/}
  \end{center}
\end{frame}

\begin{frame}
  \begin{center}
    \Huge
    Next up\ldots programming tasks!
    https://github.com/york-haskell
  \end{center}
\end{frame}

\end{document}